\section{来源审计：一堂课、两条神经科学路径}

这堂课不是单一论文讲座，而是两位讲者从不同脑区与不同计算层级，分别走向 Transformer-like computation。前 42 分钟，Trenton Bricken 从 Sparse Distributed Memory（SDM，稀疏分布式记忆）的高维几何出发，解释 hypersphere intersection 为什么产生近似 softmax 的指数权重，并提出 cerebellum（小脑）可能实现类似读写电路。后 39 分钟，Will Dorrell 从 cognitive map、hippocampus 与 entorhinal cortex 出发，介绍 Tolman--Eichenbaum Machine（TEM）如何分离结构与内容，以及 modern Hopfield memory 怎样把它改写成 attention。

\subsection{纠正旧稿的错误视频与错误范围}

旧稿链接的 \texttt{dEFn6nnoC-8} 是 Trenton Bricken 于 2025 年 3 月 31 日完成、2025 年 5 月上传的博士答辩，并非本课。正确官方录像是 Stanford Online 的 \texttt{L4DC7e6g2iI}，课堂日期为 2023 年 3 月 7 日。旧稿还只覆盖第一位讲者，并扩写出 MoE 路由、生产 dashboard、治理 playbook、incident response 和跨模态部署等课堂没有出现的内容；这些段落在本次重写中全部删除。

\begin{longtable}{@{}L{0.18\textwidth}L{0.29\textwidth}L{0.41\textwidth}@{}}
\toprule
证据层 & 本地材料 & 教学职责 \\
\midrule
官方录像 & 1:22:13、两位讲者、66 个最终教学状态 & 决定课堂顺序、公式、手写推导、问答与结论 \\
官方人工字幕 & 1,950 条 caption 与 timed transcript & 恢复动机、类比限制、实验边界和神经科学细节 \\
第一讲论文 & Attention Approximates SDM & 核对高维交集、指数近似、连续 SDM 与 learned coefficient \\
第二讲论文 & TEM 与 hippocampal Transformer 论文 & 核对结构/内容分解、记忆更新、grid-like code 与 attention 对应 \\
视觉审计 & 40 张标准 slide + 26 张手写最终状态 & 完整覆盖两段讲座，避免每一笔书写都变成重复截图 \\
\bottomrule
\end{longtable}

\begin{warningbox}{历史与证据边界}
本讲重建 2023-03-07 的课堂。数学对应说明“某种计算可以由另一种形式近似实现”，不等于证明大脑运行训练好的 Transformer；脑区映射说明“电路具备候选实现组件”，不等于已经被实验确认；TEM 与 Transformer 的相似性发生在记忆检索、位置状态与结构泛化层面，也不是逐神经元的解剖同一。
\end{warningbox}

\subsection{第一讲的研究问题与路线}

标题页给出 Bricken 与 Cengiz Pehlevan 的工作；动机页把问题压缩为一句话：Transformer attention 很强，但 softmax 是 heuristic，能否由简单高维向量性质自然产生？Agenda 页坚持先视觉直觉、后数学推导，再讨论 Transformer 组件与生物可实现性。这种顺序很重要，因为高维球面交集若直接写成特殊函数，会掩盖它作为“重叠地址数量”的记忆直觉。

\lecturefigure{slide-01-title-attention-sdm.jpg}{第一讲：Attention Approximates Sparse Distributed Memory}{官方录像恢复幻灯片 V001；00:00:22}

\lecturefigure{slide-02-why-care.jpg}{为什么值得研究：从 heuristic softmax 到高维记忆几何}{官方录像恢复幻灯片 V002；00:01:20}

\lecturefigure{slide-03-presentation-overview.jpg}{第一讲路线：SDM、attention、对应关系与 biological plausibility}{官方录像恢复幻灯片 V003；00:01:52}

\teachervoice{Bricken 明确说自己会先给 high-level visual intuition，再进入数学，并鼓励学生随时打断提问。讲义因此沿用同一顺序：先把写入、读取和交集画清楚，再推导指数权重；这不是回避公式，而是让每个符号先有操作语义。}

\subsection{两讲共同追问的三层问题}

虽然 SDM 与 TEM 来自不同理论传统，两讲都在追问三层关系：\textbf{算法层}上，模型如何检索相关记忆；\textbf{表示层}上，结构与内容怎样编码；\textbf{实现层}上，哪些神经电路或可学习矩阵能够执行计算。课程价值不在于把三层压成一句“脑就是 Transformer”，而在于比较跨层映射何时精确、何时只是启发。

\begin{importantbox}{阅读全课的证据阶梯}
依次区分：公式等价或近似、训练模型中的经验结果、神经表征相似、候选电路可实现性、真实脑内因果证据。越往后证据要求越高。一个漂亮的 attention 公式不能直接证明脑区功能，一个相似的 grid pattern 也不能直接证明两个系统使用同一算法。
\end{importantbox}

\subsection{本章小结}

正确课程源包含两位讲者和两条互补路线。第一条从高维联想记忆解释 attention 权重，第二条从认知地图与 hippocampal memory 解释结构泛化和 attention-like retrieval。所有后续内容都必须保持算法、表示与生物实现之间的证据层级。

